%%%%%%%%%%%%%%%%%%%%%%%%%%%%%%%%%%%%%%%%%%%%%%%%%%%%%%%%%%%%%%%%%%%
%                                                                 %
%                          CHAPTER FOUR                          %
%                                                                 %
%%%%%%%%%%%%%%%%%%%%%%%%%%%%%%%%%%%%%%%%%%%%%%%%%%%%%%%%%%%%%%%%%%%

\chapter{Προτεινόμενη Μέθοδος}
\label{ch:method}

Η παρούσα εργασία εντάσσεται στη γραμμή έρευνας που ανέπτυξε το πλαίσιο
FinRL~\cite{liu2018finrl} για την εφαρμογή αλγορίθμων βαθιάς ενισχυτικής μάθησης
σε χρηματοοικονομικά περιβάλλοντα, επεκτάθηκε στις αγορές κρυπτονομισμάτων
μέσω του FinRL-Crypto~\cite{gort2022finrlcrypto}, και πιο πρόσφατα εμπλουτίστηκε
με σήματα παραγόμενα από Μεγάλα Γλωσσικά Μοντέλα (LLM) στην εργασία
FinRL-DeepSeek~\cite{benhenda2025finrldeepseek}. Η παρούσα εργασία υιοθετεί αυτή
τη γενική κατεύθυνση και την επεκτείνει σε τρεις αλληλοσυμπληρωματικούς άξονες.
Πρώτον, αντικαθιστά τον πράκτορα πρόσθιας τροφοδότησης του FinRL-DeepSeek με έναν
αναδρομικό πράκτορα PPO-LSTM που χειρίζεται ρητά την κρυφή του κατάσταση,
αντιμετωπίζοντας το πρόβλημα συναλλαγών ως Μερικώς Παρατηρήσιμη Μαρκοβιανή
Διαδικασία Απόφασης (POMDP). Δεύτερον, εισάγει Περιορισμένες Μαρκοβιανές
Διαδικασίες Απόφασης (CMDP) με χαλάρωση Lagrange για την ενσωμάτωση περιορισμών
ρευστότητας και διαχείρισης κινδύνου κατά την εκπαίδευση. Τρίτον, ενσωματώνει
σήμα συναισθήματος παραγόμενο από εξειδικευμένο LLM~\cite{yang2023fingpt} σε
ωριαία δεδομένα δέκα κρυπτονομισμάτων, αξιοποιώντας το τόσο ως χαρακτηριστικό
κατάστασης όσο και ως μηχανισμό προσαρμογής της έντασης των δράσεων.

Το κεφάλαιο οργανώνεται ως εξής. Η Ενότητα~\ref{sec:method-overview} παρέχει μια
συνολική επισκόπηση της αρχιτεκτονικής και τοποθετεί την εργασία σε σχέση με τα
ανωτέρω προηγούμενα πλαίσια. Η Ενότητα~\ref{sec:method-data} παρουσιάζει τα
δεδομένα και τη διαδικασία εξαγωγής χαρακτηριστικών. Η
Ενότητα~\ref{sec:method-arch} περιγράφει την αρχιτεκτονική Actor-Critic με LSTM,
ενώ η Ενότητα~\ref{sec:method-env} διατυπώνει τυπικά το περιβάλλον συναλλαγών,
τον χώρο καταστάσεων, τον χώρο δράσεων και τις συναρτήσεις ανταμοιβής. Η
Ενότητα~\ref{sec:method-training} παρουσιάζει τον αλγόριθμο εκπαίδευσης, με τις
επιμέρους συνιστώσες της συνολικής συνάρτησης απώλειας και τον τρόπο επέκτασής
της στο πλαίσιο CMDP. Τέλος, η Ενότητα~\ref{sec:method-validation} περιγράφει τη
μεθοδολογία επικύρωσης και επιλογής μοντέλου. Η περιγραφή παραμένει σκόπιμα σε
υψηλό επίπεδο, καθώς οι ακριβείς ρυθμίσεις, τα σύνολα δεδομένων, οι
υπερπαράμετροι και τα πρωτόκολλα εκτέλεσης παρουσιάζονται αναλυτικά στο
Κεφάλαιο~\ref{ch:experiments}.


\section{Επισκόπηση της Προτεινόμενης Προσέγγισης}
\label{sec:method-overview}

Η προτεινόμενη μέθοδος αντιμετωπίζει τη διαχείριση χαρτοφυλακίου κρυπτονομισμάτων
ως πρόβλημα βαθιάς ενισχυτικής μάθησης σε συνθήκες μερικής παρατηρησιμότητας.
Ο πυρήνας του συστήματος είναι ένας πράκτορας Actor-Critic με αναδρομικό στρώμα
LSTM, ο οποίος λαμβάνει αποφάσεις κατανομής χαρτοφυλακίου σε $N = 10$
κρυπτονομίσματα σε ωριαία βάση, βελτιστοποιώντας μια συνάρτηση ανταμοιβής
βασισμένη στην ενεργό απόδοση (active return) του χαρτοφυλακίου.

\subsection*{Το Πρόβλημα: Παθητική Κατάρρευση και Μερική Παρατηρησιμότητα}

Ο κεντρικός μεθοδολογικός άξονας της εργασίας πηγάζει από μια επαναλαμβανόμενη
παθογένεια που παρατηρείται κατά την εκπαίδευση πρακτόρων DRL για συναλλαγές. Σε
συνθήκες χαμηλού λόγου σήματος προς θόρυβο, ο πράκτορας τείνει να συγκλίνει σε μια
εκφυλισμένη πολιτική χαμηλού κύκλου εργασιών, η οποία είτε παραμένει σχεδόν
πλήρως σε μετρητά είτε σχεδόν πλήρως επενδεδυμένη, αναπαράγοντας ουσιαστικά τη
συμπεριφορά ενός παθητικού δείκτη αναφοράς ισαξίας τύπου buy-and-hold. Η
συμπεριφορά αυτή, που στην παρούσα εργασία αναφέρεται ως \emph{παθητική
κατάρρευση} (passive collapse), είναι ένα τοπικό ελάχιστο ασφαλείας, ο πράκτορας
αποφεύγει το κόστος και τον κίνδυνο της ενεργού διαπραγμάτευσης χωρίς να
παράγει υπεραπόδοση. Η αντιμετώπιση της παθητικής κατάρρευσης συνιστά το νήμα που
διατρέχει και υποκινεί το σύνολο των σχεδιαστικών επιλογών που περιγράφονται στο
παρόν κεφάλαιο.

Η δυσκολία επιτείνεται από δύο εγγενή χαρακτηριστικά του προβλήματος. Πρώτον, ο
πράκτορας δεν παρατηρεί το πραγματικό υποκείμενο καθεστώς της αγοράς αλλά μόνο
ένα θορυβώδες, μερικό αποτύπωμά του, γεγονός που εντάσσει τη διαπραγμάτευση στην
κατηγορία των Μερικώς Παρατηρήσιμων Μαρκοβιανών Διαδικασιών Απόφασης (POMDP),
όπως αναλύεται στην Ενότητα~\ref{subsec:pomdp}. Δεύτερον, οι αγορές
κρυπτονομισμάτων λειτουργούν αδιάλειπτα όλο το εικοσιτετράωρο και τις επτά ημέρες
της εβδομάδας, παρουσιάζοντας εντονότερη μη στασιμότητα και απότομες μεταβολές
καθεστώτος σε σύγκριση με τις παραδοσιακές αγορές μετοχών. Σε ένα τέτοιο
περιβάλλον, ένας πράκτορας που στηρίζεται αποκλειστικά σε τιμές και τεχνικούς
δείκτες διαθέτει περιορισμένη πληροφορία για την έγκαιρη αναγνώριση των αλλαγών
καθεστώτος.

\subsection*{Η Συνεισφορά σε Τρεις Άξονες}

Σε απάντηση των ανωτέρω, η προτεινόμενη μέθοδος αρθρώνεται γύρω από τρεις
αλληλοσυμπληρωματικούς άξονες, καθένας από τους οποίους στοχεύει συγκεκριμένη
πτυχή του προβλήματος.

Πρώτον, ένας \emph{αναδρομικός πράκτορας POMDP}. Η αρχιτεκτονική Actor-Critic με
στρώμα LSTM που διατηρεί την εσωτερική του κρυφή κατάσταση αντιμετωπίζει τη
μερική παρατηρησιμότητα, χρησιμοποιώντας την κρυφή κατάσταση ως μαθημένη,
συμπυκνωμένη εκτίμηση της κατάστασης πεποίθησης (Ενότητα~\ref{subsec:pomdp-agents}).
Με αυτόν τον τρόπο ο πράκτορας αποκτά πρόσβαση σε χρονική, σχετική με το καθεστώς
πληροφορία, αντί να καταφεύγει σε μια στατική κατανομή, και η ίδια η ανταμοιβή
ενεργού απόδοσης στοχεύει ευθέως την παθητική κατάρρευση συγκρίνοντας τον
πράκτορα με τον δείκτη αναφοράς αντί με την απόλυτη απόδοση.

Δεύτερον, η \emph{διαμόρφωση ρευστότητας μέσω CMDP}. Οι απαιτήσεις ρευστότητας
και διαχείρισης κινδύνου εκφράζονται ως ρητοί περιορισμοί μέσω του πλαισίου των
Περιορισμένων Μαρκοβιανών Διαδικασιών Απόφασης (CMDP) με χαλάρωση Lagrange
(Ενότητα~\ref{subsec:cmdp}). Αντί η πολιτική να εκφυλίζεται σε μια εκφυλισμένη
γωνία ασφαλείας, ο εκπαιδεύσιμος πολλαπλασιαστής Lagrange την οδηγεί προς ένα μη
εκφυλισμένο σημείο ισορροπίας που σέβεται τους περιορισμούς ρευστότητας.

Τρίτον, η \emph{διττή αξιοποίηση σήματος συναισθήματος από LLM}. Σήμα
συναισθήματος, παραγόμενο από το εξειδικευμένο γλωσσικό μοντέλο
FinGPT~\cite{yang2023fingpt}, εισάγει εξωγενή πληροφορία πέραν των τιμών και των
τεχνικών δεικτών. Το σήμα αξιοποιείται με δύο συμπληρωματικούς τρόπους, ως
πρόσθετο χαρακτηριστικό της κατάστασης και ως μηχανισμός διαμόρφωσης της έντασης
των δράσεων.

\paragraph{Σχέση με το FinRL-DeepSeek και συνεισφορά της εργασίας.}
Η προτεινόμενη μέθοδος εκκινεί από την ιδέα του FinRL-DeepSeek~\cite{benhenda2025finrldeepseek},
το οποίο ενσωματώνει σήματα κινδύνου και συναισθήματος παραγόμενα από LLM σε έναν
πράκτορα DRL για συναλλαγές μετοχών, χτισμένο πάνω στο πλαίσιο
FinRL~\cite{liu2018finrl}. Η παρούσα εργασία επεκτείνει και διαφοροποιείται από
αυτή τη βάση στους εξής άξονες.
\begin{itemize}
    \item \textbf{Πεδίο εφαρμογής}: μεταφορά από τις παραδοσιακές αγορές μετοχών
    στις αγορές κρυπτονομισμάτων, οι οποίες λειτουργούν 24/7 και χαρακτηρίζονται
    από εντονότερη μη στασιμότητα και μεταβλητότητα.
    \item \textbf{Αρχιτεκτονική πράκτορα}: αντικατάσταση του πράκτορα πρόσθιας
    τροφοδότησης με έναν πράκτορα PPO-LSTM που αντιμετωπίζει το περιβάλλον ως POMDP.
    \item \textbf{Διαμόρφωση ρευστότητας}: επέκταση του PPO στο πλαίσιο CMDP μέσω
    χαλάρωσης Lagrange, ώστε οι περιορισμοί ρευστότητας να εκφράζονται ρητά και
    όχι ως χειροκίνητα ρυθμισμένη ποινή στην ανταμοιβή.
    \item \textbf{Υπολογιστική υποδομή και επικύρωση}: επανασχεδιασμός του αγωγού
    εκπαίδευσης για πλήρη αξιοποίηση GPU μέσω διανυσματικοποιημένων περιβαλλόντων,
    σε συνδυασμό με αυστηρή διαδικασία επικύρωσης βασισμένη σε CPCV.
\end{itemize}


\section{Δεδομένα και Εξαγωγή Χαρακτηριστικών}
\label{sec:method-data}

Τα δεδομένα που χρησιμοποιούνται για την εκπαίδευση και αξιολόγηση του πράκτορα
προέρχονται από ιστορικά ωριαία δεδομένα τύπου OHLCV (Άνοιγμα, Υψηλή, Χαμηλή,
Κλείσιμο, Όγκος) για $N = 10$ κρυπτονομίσματα: ADA, ATOM, BTC, ETC, ETH, NEO,
TRX, VET, XLM και XRP. Για κάθε περιουσιακό στοιχείο $i$ και χρονικό βήμα $t$,
τα ακολουθα πεδία αποτελούν την αρχική αναπαράσταση της αγοράς.
\begin{itemize}
    \item $p^i_o(t)$: τιμή ανοίγματος (Open).
    \item $p^i_h(t)$: υψηλότερη τιμή περιόδου (High).
    \item $p^i_l(t)$: χαμηλότερη τιμή περιόδου (Low).
    \item $p^i_c(t)$: τιμή κλεισίματος (Close).
    \item $v^i(t)$: όγκος συναλλαγών (Volume).
\end{itemize}

Η χρονική περίοδος εκπαίδευσης εκτείνεται από την 29η Απριλίου 2019 έως την 21η
Νοεμβρίου 2024, ενώ η περίοδος αξιολόγησης καλύπτει το διάστημα από την 22α
Νοεμβρίου 2024 έως την 9η Νοεμβρίου 2025. Ο διαχωρισμός αυτός διασφαλίζει ότι
η τελική μέτρηση της επίδοσης πραγματοποιείται αποκλειστικά σε δεδομένα που δεν
χρησιμοποιήθηκαν σε κανένα στάδιο της εκπαίδευσης ή της επιλογής μοντέλου.

\subsection*{Εξαγωγή Τεχνικών Χαρακτηριστικών}

Η άμεση χρήση των τιμών OHLCV ως είσοδος στον πράκτορα θα εισήγαγε σημαντικές
κλιμακικές διαφορές μεταξύ περιουσιακών στοιχείων που διαπραγματεύονται σε
εντελώς διαφορετικές τιμές (π.χ.\ BTC στα δεκάδες χιλιάδες δολάρια έναντι XRP
στα κλάσματα δολαρίου). Για τον λόγο αυτό, εφαρμόζεται επεξεργασία δεδομένων
που παράγει χαρακτηριστικά ανεξάρτητα από την απόλυτη κλίμακα των τιμών.

Το βασικό σύνολο χαρακτηριστικών (base features) περιλαμβάνει οκτώ τεχνικούς
δείκτες που υπολογίζονται ανά περιουσιακό στοιχείο και χρονικό βήμα: τη
διαφορά κινητών μέσων (MACD), τα άνω και κάτω όρια Bollinger (boll\_ub,
boll\_lb), τον δείκτη σχετικής ισχύος (RSI-30), τον δείκτη καναλιού εμπορευμάτων
(CCI-30), τον δείκτη κατεύθυνσης (DX-30), και τους απλούς κινητούς μέσους 30
και 60 βημάτων (SMA-30, SMA-60). Η εργασία διερευνά επίσης επιπλέον σύνολα
χαρακτηριστικών, που περιλαμβάνουν χαρακτηριστικά ορμής (momentum features)
και συνδυαστικά προεπιλεγμένα σύνολα που βασίζονται σε παράθυρα 24 ωρών
(combo presets A έως E), τα οποία παρουσιάζονται αναλυτικά στο
Κεφάλαιο~\ref{ch:experiments}.

Για κάθε τεχνικό χαρακτηριστικό $f$ του περιουσιακού στοιχείου $i$
εφαρμόζεται τυποποίηση z-score με χρήση των στατιστικών της περιόδου
εκπαίδευσης:
\begin{equation}
    \tilde{f}^i(t) = \frac{f^i(t) - \mu^i_f}{\sigma^i_f}
    \label{eq:zscore}
\end{equation}
όπου $\mu^i_f$ και $\sigma^i_f$ είναι ο μέσος και η τυπική απόκλιση του
χαρακτηριστικού $f$ για το περιουσιακό στοιχείο $i$ επί της περιόδου
εκπαίδευσης. Οι τυποποιημένες τιμές περιορίζονται στο διάστημα $[-6, 6]$
για την αποφυγή ακραίων τιμών.

\subsection*{Ενσωμάτωση Σήματος Συναισθήματος}

Πέραν των τεχνικών χαρακτηριστικών, η εργασία διερευνά την ενσωμάτωση σήματος
συναισθήματος ως φορέα εξωγενούς πληροφορίας, πέραν εκείνης που εμπεριέχεται στις
τιμές και στους τεχνικούς δείκτες. Το σήμα παράγεται από το εξειδικευμένο
γλωσσικό μοντέλο FinGPT~\cite{yang2023fingpt}, το οποίο μικρο-προσαρμόζεται στη
χρηματοοικονομική γλώσσα μέσω Low-Rank Adaptation (LoRA)~\cite{hu2021lora}, όπως
περιγράφεται στην Ενότητα~\ref{subsec:finetuning}. Η προσέγγιση αυτή επιτρέπει
αποδοτική εξαγωγή τιμών συναισθήματος από ειδήσεις και αναρτήσεις κοινωνικής
δικτύωσης σχετικά με τα εξεταζόμενα κρυπτονομίσματα.

Το παραγόμενο αρχικό σήμα είναι μια διακριτή βαθμίδα συναισθήματος σε ακέραια
κλίμακα πέντε επιπέδων (από έντονα αρνητικό έως έντονα θετικό). Η βαθμίδα αυτή
κωδικοποιείται ως δυαδικό διάνυσμα τύπου one-hot πέντε στοιχείων, όσες και οι
κλάσεις της κλίμακας, ανά περιουσιακό στοιχείο και χρονικό βήμα. Το
κωδικοποιημένο σήμα είτε ενσωματώνεται στο διάνυσμα κατάστασης είτε
χρησιμοποιείται για τη διαμόρφωση της έντασης των δράσεων, όπως περιγράφεται στην
Ενότητα~\ref{sec:method-env}. Ο ακριβής συγχρονισμός του σήματος με τα χρονικά
βήματα της αγοράς είναι κρίσιμος, ώστε σε κάθε βήμα να αξιοποιείται μόνο
πληροφορία διαθέσιμη μέχρι τη στιγμή εκείνη και να αποφεύγεται η διαρροή
μελλοντικής πληροφορίας (look-ahead bias).


\section{Αρχιτεκτονική του Πράκτορα}
\label{sec:method-arch}

Η επιλογή αναδρομικής αρχιτεκτονικής αποτελεί άμεση απόκριση στη μερική
παρατηρησιμότητα του προβλήματος. Όπως αναλύεται στην
Ενότητα~\ref{subsec:pomdp-agents}, η κρυφή κατάσταση ενός αναδρομικού δικτύου
λειτουργεί ως μαθημένη, πεπερασμένης διάστασης εκτίμηση της κατάστασης πεποίθησης
της POMDP, συμπυκνώνοντας το σχετικό ιστορικό παρατηρήσεων χωρίς ρητή γνώση του
μοντέλου παρατήρησης. Παρέχοντας στον πράκτορα πρόσβαση σε αυτή τη χρονική
αναπαράσταση του καθεστώτος, η αρχιτεκτονική του επιτρέπει να διαφοροποιεί τη
συμπεριφορά του ανά καθεστώς αγοράς, αντί να καταφεύγει στη στατική κατανομή που
χαρακτηρίζει την παθητική κατάρρευση.

Ο πράκτορας υλοποιεί το παράδειγμα Actor-Critic με κοινό αναδρομικό στρώμα
LSTM, του οποίου η θεωρητική θεμελίωση παρουσιάζεται στην
Ενότητα~\ref{subsec:lstm}. Σε αντίθεση με αρχιτεκτονικές που τροφοδοτούν ένα
σταθερό παράθυρο παρελθόντων παρατηρήσεων (sliding window) ως είσοδο, η παρούσα
προσέγγιση υιοθετεί αρχιτεκτονική LSTM με διατήρηση της εσωτερικής αναδρομικής
κατάστασης $(h_t, c_t)$ μεταξύ διαδοχικών χρονικών βημάτων. Σε κάθε βήμα $t$, ο
πράκτορας λαμβάνει μία και μόνο τρέχουσα παρατήρηση $o_t \in \mathbb{R}^{d}$, ενώ
η χρονική πληροφορία συσσωρεύεται εσωτερικά μέσω της αναδρομικής κατάστασης,
χωρίς να απαιτείται η διατήρηση ιστορικού στον χώρο παρατήρησης. Η προσέγγιση του
παραθύρου παρελθόντων παρατηρήσεων αποτελεί μια εξίσου θεμιτή αλλά διαφορετική
σχεδιαστική επιλογή, η οποία ωστόσο μεταθέτει το βάρος της χρονικής μνήμης στον
χώρο εισόδου και αυξάνει τη διάστασή του, αντί να την κωδικοποιεί στην κρυφή
κατάσταση του δικτύου.

\subsection*{Δομή Δικτύου}

Η αρχιτεκτονική αποτελείται από τα εξής επίπεδα, κοινά για τον Actor και τον
Critic:
\begin{enumerate}
    \item \textbf{Στρώμα LSTM}: δέχεται την παρατήρηση $o_t$ ενός βήματος
    $\bigl(\text{διαστάσεων } (1, d)\bigr)$ και ενημερώνει την κρυφή κατάσταση
    $(h_t, c_t)$, παράγοντας έξοδο $\ell_t \in \mathbb{R}^{d_h}$ όπου $d_h$
    είναι ο αριθμός των κρυφών μονάδων.
    \item \textbf{Βάση MLP}: τροφοδοτείται με την έξοδο $\ell_t$ του LSTM και
    παράγει μια αναπαράσταση πολιτικής/αξίας μέσω ενός ή δύο γραμμικών
    επιπέδων με συνάρτηση ενεργοποίησης.
    \item \textbf{Κεφαλή Actor}: γραμμικό επίπεδο που παράγει τις παραμέτρους
    $(\mu_t, \log \sigma_t)$ μιας γκαουσιανής κατανομής επί του χώρου δράσεων
    $\mathcal{A} = [-1, 1]^{N}$.
    \item \textbf{Κεφαλή Critic}: γραμμικό επίπεδο που παράγει εκτίμηση της
    συνάρτησης αξίας $V_\phi(o_t, h_t)$.
\end{enumerate}

Η ρητή διαχείριση της κρυφής κατάστασης αποτελεί κρίσιμη συνιστώσα ορθότητας,
καθώς η κατάσταση $(h_t, c_t)$ μηδενίζεται στα όρια των επεισοδίων και πριν από
κάθε πέρασμα ενημέρωσης PPO, αποφεύγοντας τη διαρροή πληροφορίας μεταξύ
ασύνδετων τροχιών.

\subsection*{Παράλληλη Εκπαίδευση με Διανυσματικοποιημένα Περιβάλλοντα}

Για την αποδοτική συλλογή εμπειρίας, χρησιμοποιείται διανυσματικοποιημένο
περιβάλλον τύπου \texttt{SyncVectorEnv}, το οποίο εκτελεί $M$ παράλληλα
αντίγραφα του περιβάλλοντος ταυτόχρονα. Σε κάθε βήμα, ο πράκτορας επεξεργάζεται
ένα δέσμη παρατηρήσεων $O_t \in \mathbb{R}^{M \times d}$ και παράγει αντίστοιχες
δράσεις για όλα τα περιβάλλοντα, αξιοποιώντας τον παραλληλισμό της GPU. Η
κρυφή κατάσταση διατηρείται ξεχωριστά για κάθε παράλληλο περιβάλλον. Ο
αγωγός εκπαίδευσης εκτελεί την αλληλεπίδραση πράκτορα-περιβάλλοντος στην
CPU και μεταφέρει τμηματικά τα δεδομένα στη GPU κατά τις ενημερώσεις PPO,
ενώ η αποκομμένη οπισθοδιάδοση στον χρόνο (truncated backpropagation through
time, BPTT) περιορίζει τις απαιτήσεις μνήμης GPU σε κατανεκτά επίπεδα.


\section{Περιβάλλον Συναλλαγών}
\label{sec:method-env}

Το περιβάλλον μοντελοποιεί τη διαχείριση χαρτοφυλακίου $N$ κρυπτονομισμάτων
ως Μερικώς Παρατηρήσιμη Μαρκοβιανή Διαδικασία Απόφασης
(POMDP)~\cite{kaelbling1998pomdp}, σύμφωνα με το πλαίσιο της
Ενότητας~\ref{subsec:pomdp}. Ο πράκτορας δεν έχει πρόσβαση στο υποκείμενο
καθεστώς της αγοράς, αλλά μόνο σε ένα θορυβώδες αποτύπωμά του μέσω
παρατηρήσεων ωριαίων δεδομένων.

\subsection*{Χώρος Καταστάσεων}

Σε κάθε χρονικό βήμα $t$, η παρατήρηση $o_t$ συγκροτείται από τρεις
συνιστώσες:
\begin{equation}
    o_t = \bigl[ w^{\text{cash}}_t, \; \mathbf{w}^{\text{hold}}_t, \;
    \mathbf{x}^{\text{tech}}_t \bigr] \in \mathbb{R}^{d},
    \label{eq:obs}
\end{equation}
όπου:
\begin{itemize}
    \item $w^{\text{cash}}_t \in [0, 1]$ είναι ο λόγος των διαθέσιμων
    μετρητών προς τη συνολική αξία του χαρτοφυλακίου (cash ratio).
    \item $\mathbf{w}^{\text{hold}}_t \in \mathbb{R}^{N}$ είναι το διάνυσμα
    των τρεχουσών συμμετοχών (holdings) ανά περιουσιακό στοιχείο, εκφρασμένο
    ως κλάσμα της συνολικής αξίας χαρτοφυλακίου.
    \item $\mathbf{x}^{\text{tech}}_t \in \mathbb{R}^{N \cdot k}$ είναι το
    διάνυσμα τεχνικών χαρακτηριστικών για όλα τα $N$ περιουσιακά στοιχεία,
    με $k$ χαρακτηριστικά ανά στοιχείο.
\end{itemize}
Στη διαμόρφωση με σήμα συναισθήματος, στο διάνυσμα $o_t$ προστίθεται και το
αντίστοιχο one-hot σήμα $\mathbf{s}_t$, το οποίο κωδικοποιεί τη διακριτή βαθμίδα
συναισθήματος πέντε επιπέδων, αυξάνοντας αντιστοίχως τη συνολική διάσταση $d$
της παρατήρησης.

\subsection*{Χώρος Δράσεων}

Ο χώρος δράσεων είναι συνεχής:
\begin{equation}
    \mathcal{A} = [-1, 1]^{N}.
    \label{eq:action-space}
\end{equation}
Κάθε στοιχείο $a^i_t$ της δράσης ερμηνεύεται ως το \emph{στοχευόμενο βάρος}
(target weight) του $i$-οστού περιουσιακού στοιχείου στο χαρτοφυλάκιο, δηλαδή ως
το επιθυμητό κλάσμα της συνολικής αξίας που πρέπει να επενδυθεί σε αυτό. Σε κάθε
βήμα το περιβάλλον προβάλλει το διάνυσμα δράσεων στον εφικτό χώρο, επιβάλλοντας
ότι το άθροισμα των βαρών δεν υπερβαίνει τη μονάδα (αποκλεισμός έμμεσης
μόχλευσης), και ακολούθως εκτελεί τις απαιτούμενες αγορές και πωλήσεις ώστε η
τρέχουσα σύνθεση να συγκλίνει προς το στόχο, αφαιρώντας το αντίστοιχο κόστος
συναλλαγής. Έτσι, η ίδια η δράση καθορίζει την κατανομή στόχου και όχι μια
αυξητική μεταβολή της τρέχουσας θέσης.

Οι δράσεις δειγματοληπτούνται από μια γκαουσιανή πολιτική
$\pi_\theta(a \mid o_t, h_t) = \mathcal{N}(\mu_\theta, \sigma_\theta^2)$. Το
ζεύγος $\bigl(a_t, \log \pi_\theta(a_t \mid o_t, h_t)\bigr)$ αποθηκεύεται πριν
από την προβολή στον εφικτό χώρο, ώστε ο λόγος πιθανοτήτων του PPO να υπολογίζεται
ως προς τη δράση που πράγματι δειγματολήφθηκε από την πολιτική και να διατηρείται
η ορθότητα της ενημέρωσης.

\subsection*{Αρχική Κατάσταση και Εργοδικότητα της Πολιτικής}

Ένα λεπτό αλλά ουσιώδες σημείο σχεδιασμού αφορά την αρχικοποίηση του
χαρτοφυλακίου στην αρχή κάθε επεισοδίου. Επειδή η δράση αντιστοιχεί απευθείας σε
στοχευόμενα βάρη, μια βέλτιστη πολιτική σε ένα συνεχιζόμενο πρόβλημα επάγει μια
εργοδική Μαρκοβιανή αλυσίδα της οποίας η στάσιμη κατανομή είναι ανεξάρτητη από
την αρχική κατάσταση~\cite{sutton2018rl}, ο πράκτορας απλώς αναπροσαρμόζει τη
σύνθεση προς το προτιμώμενο σημείο ισορροπίας του, με μόνη μόνιμη επίδραση της
αρχικής κατάστασης ένα εφάπαξ μεταβατικό κόστος κύκλου εργασιών. Η παρατήρηση
αυτή έχει δύο συνέπειες. Αφενός, η αξιολόγηση μπορεί να εκκινεί από μια σταθερή,
ισορροπημένη αρχική κατάσταση (ισομερής κατανομή μεταξύ μετρητών και θέσεων),
χωρίς απώλεια γενικότητας. Αφετέρου, και κρισιμότερα για το πλαίσιο CMDP, κατά
την εκπαίδευση η αρχική κατάσταση τυχαιοποιείται ομοιόμορφα ως προς το ποσοστό
των μετρητών και την κατανομή στις θέσεις. Η τυχαιοποίηση αυτή διασφαλίζει ότι ο
πράκτορας και ο εκτιμητής αξίας του εκτίθενται σε ολόκληρο το εύρος των πιθανών
συνθέσεων χαρτοφυλακίου, συμπεριλαμβανομένων των περιοχών που παραβιάζουν τους
περιορισμούς ρευστότητας. Η κάλυψη αυτών των περιοχών είναι αναγκαία ώστε ο
εκτιμητής αξίας να μάθει την τοπολογία των ποινών και ο πολλαπλασιαστής Lagrange
να προσαρμοστεί σωστά, σύμφωνα με τα ευρήματα της βιβλιογραφίας για την ασφαλή
εξερεύνηση~\cite{garcia2015safe,tessler2018rcpo}.

\subsection*{Συνάρτηση Ανταμοιβής}

Η σχεδίαση της ανταμοιβής αποτελεί την πιο άμεση απόκριση στην παθητική
κατάρρευση. Μια ανταμοιβή που στηρίζεται στην απόλυτη απόδοση του χαρτοφυλακίου
ανταμείβει τον πράκτορα ακόμη και όταν αυτός απλώς παρακολουθεί την άνοδο της
αγοράς, αφήνοντας την παθητική στρατηγική buy-and-hold ως ελκυστικό σημείο
ισορροπίας. Αντιθέτως, η βασική συνάρτηση ανταμοιβής της εργασίας στηρίζεται στην
\emph{ενεργό απόδοση} (active return) του χαρτοφυλακίου σε σχέση με έναν δείκτη
αναφοράς ισαξίας (equal-weight buy-and-hold), έννοια συγγενή του ενεργού κέρδους
(alpha) που ορίζεται στην Ενότητα~\ref{subsec:alpha}. Με αυτόν τον τρόπο, η
απλή αναπαραγωγή της συμπεριφοράς του δείκτη αναφοράς αποδίδει μηδενική ανταμοιβή
και ο πράκτορας ανταμείβεται μόνο για την υπεραπόδοση έναντι αυτού.

Σε κάθε βήμα $t$, η ανταμοιβή ορίζεται ως η διαφορά της λογαριθμικής απόδοσης του
χαρτοφυλακίου από εκείνη του δείκτη αναφοράς:
\begin{equation}
    r_t = R^{\text{port}}_t - R^{\text{bench}}_t,
    \label{eq:reward}
\end{equation}
όπου
\begin{equation}
    R^{\text{port}}_t = \ln\!\frac{A_t}{A_{t-1}}, \qquad
    R^{\text{bench}}_t = \ln\!\left(\frac{1}{N}\sum_{i=1}^{N}
    \frac{p^i_c(t)}{p^i_c(t-1)}\right),
    \label{eq:reward-components}
\end{equation}
με $A_t$ τη συνολική αξία του χαρτοφυλακίου στο βήμα $t$ και $p^i_c(t)$ την τιμή
κλεισίματος του $i$-οστού στοιχείου. Το κόστος συναλλαγής δεν εισέρχεται ως
ξεχωριστός όρος ποινής στην ανταμοιβή, αλλά αφαιρείται απευθείας από την αξία του
χαρτοφυλακίου κατά την εκτέλεση κάθε αγοράς ή πώλησης, οπότε αντανακλάται φυσικά
στον όρο $R^{\text{port}}_t$. Έτσι, ένας υψηλός κύκλος εργασιών τιμωρείται
ενδογενώς μέσω του σωρευτικού κόστους συναλλαγών, χωρίς ανάγκη χειροκίνητης
στάθμισης ενός όρου ποινής. Το ποσοστιαίο κόστος αγοράς και πώλησης ορίζεται ως
ρυθμιζόμενη παράμετρος, της οποίας οι συγκεκριμένες τιμές που εξετάστηκαν
παρουσιάζονται στο Κεφάλαιο~\ref{ch:experiments}. Η αρχική κεφαλαιοποίηση
ορίζεται σε 1.000.000 νομισματικές μονάδες.

Η εργασία διερευνά επίσης παραλλαγές της ανταμοιβής, που περιλαμβάνουν
συνάρτηση βασισμένη στον Λόγο Sharpe (Ενότητα~\ref{subsec:alpha}), ανταμοιβή $N$
βημάτων (N-step reward) και παραλλαγές με ρητή ποινή κύκλου εργασιών, τα
αποτελέσματα των οποίων αξιολογούνται στο Κεφάλαιο~\ref{ch:experiments}.

\subsection*{Περιορισμοί Ρευστότητας μέσω CMDP}

Η ανταμοιβή ενεργού απόδοσης αποτρέπει την παρακολούθηση της αγοράς, αλλά δεν
εμποδίζει από μόνη της τον πράκτορα να εκφυλιστεί προς μια διαφορετική γωνία
ασφαλείας, για παράδειγμα να διατηρεί μονίμως ένα πολύ υψηλό ποσοστό μετρητών,
αποφεύγοντας κάθε κίνδυνο. Για να ωθηθεί ο πράκτορας προς ένα μη εκφυλισμένο,
ρευστό σημείο ισορροπίας, οι απαιτήσεις ρευστότητας και διαχείρισης κινδύνου
εκφράζονται ως ρητοί περιορισμοί μέσω του πλαισίου των Περιορισμένων Μαρκοβιανών
Διαδικασιών Απόφασης (CMDP)~\cite{altman1999cmdp}, του οποίου η θεωρητική
θεμελίωση παρουσιάζεται στην Ενότητα~\ref{subsec:cmdp}. Σε αρμονία με την
Εξίσωση~\eqref{eq:cmdp} του Κεφαλαίου~\ref{ch:theory}, το βελτιστοποιητικό
πρόβλημα διατυπώνεται ως:
\begin{equation}
    \max_{\theta} \; \mathbb{E}_{\pi_\theta}\!\left[\sum_{t=0}^{T} \gamma^t r_t\right]
    \quad \text{υπό τον περιορισμό} \quad
    \mathbb{E}_{\pi_\theta}\!\left[\sum_{t=0}^{T} g_t\right] \leq d_0,
    \label{eq:cmdp-method}
\end{equation}
όπου $g_t$ είναι το κόστος ρευστότητας και $d_0$ το ανεκτό όριο. Η εργασία
εξετάζει διαφορετικές διατυπώσεις του κόστους $g_t$, μεταξύ των οποίων η
παραβίαση ενός κατωφλίου ελάχιστων μετρητών (liquidity buffer), η απόκλιση από
μια ζώνη επιτρεπτών μετρητών με άνω και κάτω όριο (cash band), και μια
παραλλαγή που επιβάλλει ενεργή ρευστοποίηση θέσεων (liquify) όταν τα διαθέσιμα
μετρητά πέφτουν κάτω από ένα στόχο διαθεσίμων. Μέσω χαλάρωσης Lagrange, το
πρόβλημα μετασχηματίζεται στο εξής πρόβλημα ελαχιστομεγιστοποίησης:
\begin{equation}
    \max_{\theta} \min_{\lambda \geq 0} \;
    \mathbb{E}_{\pi_\theta}\!\left[\sum_{t=0}^{T} \gamma^t \bigl(r_t - \lambda \, g_t\bigr)\right],
    \label{eq:lagrangian}
\end{equation}
όπου $\lambda \geq 0$ είναι ο εκπαιδεύσιμος πολλαπλασιαστής Lagrange. Στην
υλοποίηση, η μη αρνητικότητα του πολλαπλασιαστή επιβάλλεται εκφράζοντάς τον ως
$\lambda = \mathrm{softplus}(\ell)$, όπου $\ell$ είναι η ελεύθερη παράμετρος που
βελτιστοποιείται, ώστε ο $\lambda$ να παραμένει αυστηρά μη αρνητικός χωρίς
ανάγκη ρητής προβολής. Κατά την εκπαίδευση, ο $\lambda$ αυξάνεται όταν ο μέσος
περιορισμός παραβιάζεται και μειώνεται όταν ικανοποιείται, επιτρέποντας στον
πράκτορα να ισορροπήσει βαθμιαία μεταξύ μεγιστοποίησης της απόδοσης και τήρησης
του περιορισμού ρευστότητας, χωρίς χειροκίνητη ρύθμιση ενός σταθερού βάρους
ποινής. Σε σύγκριση με την ενσωμάτωση του κινδύνου ως σταθερής ποινής στην
ανταμοιβή, η προσέγγιση αυτή οδηγεί σε πιο ερμηνεύσιμη και ελεγχόμενη
συμπεριφορά~\cite{tessler2018rcpo,achiam2017cpo}.

\subsection*{Διττή Αξιοποίηση του Σήματος Συναισθήματος}

Το σήμα συναισθήματος αξιοποιείται με δύο συμπληρωματικούς τρόπους, σε αρμονία με
την ανάλυση της ενσωμάτωσης LLM στην ανάλυση χρηματοοικονομικού συναισθήματος
(Ενότητα~\ref{subsec:llm-sentiment}). Πρώτον, ως \emph{χαρακτηριστικό
κατάστασης}, όπου το one-hot διάνυσμα συναισθήματος $\mathbf{s}_t$ εντάσσεται
απευθείας στην παρατήρηση $o_t$, επιτρέποντας στο LSTM να ενσωματώσει την
πληροφορία του κλίματος στην εκτίμηση πεποίθησής του και να μάθει χρονικά μοτίβα
που συνδέουν τη μεταβολή του επενδυτικού κλίματος με μελλοντικές κινήσεις της
αγοράς.

Δεύτερον, ως \emph{μηχανισμός προσαρμογής της έντασης δράσεων}, όπου το σήμα
διαμορφώνει την επιθετικότητα της στρατηγικής χωρίς να μεταβάλλει την κατεύθυνσή
της. Συγκεκριμένα, ανάλογα με τη βαθμίδα του συναισθήματος, οι δράσεις αγοράς και
πώλησης πολλαπλασιάζονται με συντελεστές μεγαλύτερους ή μικρότερους της μονάδας.
Σε θετικό κλίμα ενισχύονται οι αγορές και αποδυναμώνονται οι πωλήσεις, ενώ σε
αρνητικό κλίμα ισχύει το αντίστροφο, με ουδέτερη βαθμίδα να αφήνει τις δράσεις
αμετάβλητες:
\begin{equation}
    \tilde{a}^i_t = \kappa(s_t,\, \mathrm{sgn}(a^i_t)) \cdot a^i_t,
    \label{eq:sentiment-scaling}
\end{equation}
όπου ο συντελεστής $\kappa(\cdot) > 0$ εξαρτάται από τη βαθμίδα συναισθήματος
$s_t$ και από το πρόσημο της δράσης (αγορά ή πώληση). Η προσαρμογή εφαρμόζεται
πριν από την εκτέλεση στο περιβάλλον, ενώ ο PPO αποθηκεύει το αρχικό ζεύγος
$(a^i_t, \log\pi_\theta)$ της δειγματοληφθείσας δράσης για τον ορθό υπολογισμό του
λόγου πιθανοτήτων.


\section{Αλγόριθμος Εκπαίδευσης}
\label{sec:method-training}

Η εκπαίδευση βασίζεται στον αλγόριθμο Proximal Policy Optimization
(PPO)~\cite{schulman2017ppo}, του οποίου η πλήρης διατύπωση, η συνάρτηση
αποκομμένου υποκατάστατου στόχου (CLIP) και οι επιμέρους όροι απώλειας
(πολιτικής, αξίας και εντροπίας) αναπτύσσονται στην Ενότητα~\ref{subsec:ppo}.
Η παρούσα ενότητα επικεντρώνεται στις προσαρμογές που είναι ειδικές για την
προτεινόμενη μέθοδο, δηλαδή στον χειρισμό της αναδρομικής κατάστασης κατά την
ενημέρωση, στην τυποποίηση των πλεονεκτημάτων και στην επέκταση CMDP με τον
συνεκπαιδευόμενο πολλαπλασιαστή Lagrange. Ο αλγόριθμος εκτελείται σε εποχές
(epochs), όπου σε κάθε εποχή συλλέγεται δέσμη εμπειρίας μέσω των $M$ παράλληλων
περιβαλλόντων και ακολουθούν $K$ βήματα ενημέρωσης των παραμέτρων.

\subsection*{Εκτίμηση Πλεονεκτήματος με Αναδρομικό Εκτιμητή Αξίας}

Το πλεονέκτημα $A_t$ εκτιμάται με τη Γενικευμένη Εκτίμηση Πλεονεκτήματος (GAE),
όπως ορίζεται στην Εξίσωση~\eqref{eq:gae_eq} του Κεφαλαίου~\ref{ch:theory}. Η
ιδιαιτερότητα στην παρούσα μέθοδο είναι ότι ο εκτιμητής αξίας είναι αναδρομικός,
οπότε εξαρτάται τόσο από την τρέχουσα παρατήρηση όσο και από την κρυφή κατάσταση,
$V_\phi(o_t, h_t)$. Το χρονικό σφάλμα ορίζεται αντιστοίχως ως
\begin{equation}
    \delta_t = r_t + \gamma\, V_\phi(o_{t+1}, h_{t+1}) - V_\phi(o_t, h_t),
    \label{eq:td-error}
\end{equation}
και τα προκύπτοντα πλεονεκτήματα τυποποιούνται (αφαίρεση μέσου, διαίρεση με την
τυπική απόκλιση) πριν από κάθε ενημέρωση, ώστε να σταθεροποιείται η κλίμακα της
κλίσης. Η τυποποίηση γίνεται καθολικά, επί του συνόλου της δέσμης που συλλέγεται
από όλα τα παράλληλα περιβάλλοντα, ώστε όλες οι τροχιές να συνεισφέρουν σε κοινή
κλίμακα πλεονεκτήματος.

\subsection*{Επέκταση CMDP και Ενημέρωση του Πολλαπλασιαστή}

Στη διαμόρφωση CMDP, η ανταμοιβή $r_t$ που τροφοδοτεί τον υπολογισμό του
πλεονεκτήματος αντικαθίσταται από την \emph{αυξημένη ανταμοιβή}
\begin{equation}
    r^{\text{aug}}_t = r_t - \lambda\, g_t,
    \label{eq:aug-reward}
\end{equation}
όπου $g_t$ είναι το κόστος ρευστότητας και $\lambda = \mathrm{softplus}(\ell)$ ο
μη αρνητικός πολλαπλασιαστής. Η ελεύθερη παράμετρος $\ell$ ενημερώνεται με
ανάβαση κλίσης ως προς το μέσο πλεόνασμα του περιορισμού, αυξάνοντας την ποινή
όταν ο μέσος περιορισμός παραβιάζεται και μειώνοντάς την όταν ικανοποιείται:
\begin{equation}
    \ell \leftarrow \ell + \eta_\lambda \,
    \mathrm{softplus}'(\ell)\,\bigl(\mathbb{E}[g_t] - d_0\bigr),
    \label{eq:lambda-update}
\end{equation}
με $\eta_\lambda$ τον ρυθμό μάθησης του πολλαπλασιαστή. Η συνεκπαίδευση των
παραμέτρων της πολιτικής και του $\ell$ υλοποιεί την ελαχιστομεγιστοποίηση της
Εξίσωσης~\eqref{eq:lagrangian}, οδηγώντας τον πράκτορα προς μια πολιτική που
ισορροπεί τη μεγιστοποίηση της ενεργού απόδοσης με την τήρηση των περιορισμών
ρευστότητας.

\subsection*{Αποκομμένη Οπισθοδιάδοση στον Χρόνο}

Για τον περιορισμό της μνήμης GPU κατά την εκπαίδευση LSTM, εφαρμόζεται
αποκομμένη οπισθοδιάδοση στον χρόνο (truncated BPTT) με παράθυρο $B$
βημάτων. Η τροχιά κάθε επεισοδίου τεμαχίζεται σε τμήματα μήκους $B$,
και η οπισθοδιάδοση εφαρμόζεται ανεξάρτητα σε κάθε τμήμα, με αποκοπή
της κλίσης στα όρια των τμημάτων. Η κρυφή κατάσταση μεταφέρεται μεταξύ
τμημάτων ως εκκίνηση, αλλά χωρίς διάδοση κλίσης, ισορροπώντας μεταξύ
υπολογιστικής αποδοτικότητας και ικανότητας εκμάθησης μακροπρόθεσμων
εξαρτήσεων.


\section{Μεθοδολογία Επικύρωσης και Επιλογής Μοντέλου}
\label{sec:method-validation}

Οι πράκτορες DRL σε χρηματοοικονομικά δεδομένα είναι ιδιαιτέρως ευάλωτοι στην
υπερπροσαρμογή στο παρελθόν (backtest overfitting), ζήτημα που αναλύεται στην
Ενότητα~\ref{subsec:ai-trading} και οξύνεται από τον μεγάλο χώρο υπερπαραμέτρων
και τον χαμηλό λόγο σήματος προς θόρυβο των αγορών. Για τον λόγο αυτό, η επιλογή
μοντέλου δεν στηρίζεται σε έναν απλό διαχωρισμό εκπαίδευσης και επικύρωσης, ο
οποίος παράγει μία μόνο εκτίμηση επίδοσης και είναι εύκολο να υπερβελτιστοποιηθεί.
Αντ' αυτού, υιοθετείται η μέθοδος Combinatorial Purged Cross-Validation
(CPCV)~\cite{lopezdeprado2018}, η οποία σχεδιάστηκε ειδικά για χρονοσειρές.

Η CPCV διαμερίζει την περίοδο εκπαίδευσης σε διαδοχικές ομάδες και σχηματίζει
πολλαπλούς συνδυασμούς ζωνών εκπαίδευσης και επικύρωσης. Σε κάθε συνδυασμό
εφαρμόζονται δύο μηχανισμοί προστασίας έναντι της διαρροής πληροφορίας. Η
απομόνωση (purging) αφαιρεί από τη ζώνη εκπαίδευσης τις παρατηρήσεις των οποίων η
πληροφορία επικαλύπτεται χρονικά με τη ζώνη επικύρωσης, ενώ το εμπάργκο (embargo)
εισάγει ένα χρονικό κενό μετά από κάθε ζώνη επικύρωσης, ώστε η αυτοσυσχέτιση των
χρηματοοικονομικών σειρών να μην μεταφέρει πληροφορία από το μέλλον προς το
παρελθόν. Η προσέγγιση αυτή ακολουθεί τη γραμμή εργασιών που εφαρμόζουν αυστηρή
διασταυρούμενη επικύρωση σε πράκτορες DRL για
κρυπτονομίσματα~\cite{gort2022finrlcrypto}.

Η αναζήτηση υπερπαραμέτρων διεξάγεται μέσω του πλαισίου Optuna~\cite{akiba2019optuna},
ολοκληρωτικά εντός της διαδικασίας CPCV, ώστε καμία απόφαση σχεδιασμού να μην
αξιολογηθεί έναντι του τελικού συνόλου αξιολόγησης. Για κάθε υποψήφια
παραμετροποίηση, ο πράκτορας εκπαιδεύεται και αξιολογείται σε όλους τους
συνδυασμούς ζωνών που ορίζει η CPCV, και η τελική βαθμολογία υπολογίζεται ως
μέσος όρος των επιδόσεων στις ζώνες επικύρωσης. Η άντληση της βαθμολογίας από
πολλαπλά, αλληλεπικαλυπτόμενα μονοπάτια εκπαίδευσης και επικύρωσης παρέχει μια πιο
εύρωστη εκτίμηση της γενίκευσης από έναν μεμονωμένο διαχωρισμό και μειώνει την
πιθανότητα επιλογής μιας παραμετροποίησης που αποδίδει καλά τυχαία σε μία μόνο
χρονική περίοδο.

Καθοριστικής σημασίας για την εγκυρότητα της αξιολόγησης είναι ο πλήρης
αποκλεισμός του τελικού συνόλου από κάθε στάδιο εκπαίδευσης ή επιλογής μοντέλου.
Το τελικό σύνολο αξιολόγησης, που καλύπτει την περίοδο από 22/11/2024 έως
09/11/2025, παραμένει αυστηρά απομονωμένο καθ' όλη τη διάρκεια της αναζήτησης
υπερπαραμέτρων και χρησιμοποιείται μία και μόνη φορά, στο τέλος, για την τελική
και αμερόληπτη μέτρηση της επίδοσης του επιλεγμένου μοντέλου. Στο τελικό αυτό
σύνολο, η αξιολόγηση εκκινεί από την ισορροπημένη αρχική κατάσταση που
περιγράφηκε στην Ενότητα~\ref{sec:method-env}. Το πλήρες πειραματικό πρωτόκολλο,
τα σημεία αναφοράς (baselines) και τα αποτελέσματα παρουσιάζονται στο
Κεφάλαιο~\ref{ch:experiments}.

%% --- ΤΕΛΟΣ ΚΕΦΑΛΑΙΟΥ 4 ---
